\problema{Rotting Oranges}{994}{src/01-grafos/994-rotting-oranges.cpp}{M}

Tenemos una cuadrícula donde \texttt{0} es una celda vacía, \texttt{1} es una
naranja fresca y \texttt{2} una naranja podrida. Cada minuto, toda naranja
fresca adyacente (arriba, abajo, izquierda, derecha) a una podrida se pudre.
Queremos el número de minutos hasta que no quede ninguna fresca, o
\texttt{-1} si alguna nunca se pudre.

\paragraph{Intuición.} La podredumbre se expande como una onda desde todas
las naranjas podridas \emph{al mismo tiempo}. Eso es exactamente un BFS
\emph{multi-fuente}: metemos todas las podridas en la cola y expandimos por
capas; cada capa es un minuto. Contamos las frescas al inicio y las restamos
conforme se pudren; si al final queda alguna, es imposible.

\begin{center}
\ttfamily
minuto 0\quad
\begin{tabular}{ccc}2&1&1\\1&1&0\\0&1&1\end{tabular}
\;$\rightarrow$\;
minuto 4\quad
\begin{tabular}{ccc}2&2&2\\2&2&0\\0&2&2\end{tabular}
\end{center}

\paragraph{Solución.} BFS por niveles desde todas las podridas. La clave del
manejo de casos borde: si no hay frescas al inicio, la respuesta es
\texttt{0}; si terminan las capas y aún quedan frescas, es \texttt{-1}.

\lstinputlisting{src/01-grafos/994-rotting-oranges.cpp}

\complejidad{$O(m\cdot n)$}{$O(m\cdot n)$}

\paragraph{Notas de entrevista (Amazon).} El error típico es empezar
\texttt{minutos} en $-1$ o contar un minuto de más cuando la última capa no
pudre nada. Aquí sólo incrementamos \texttt{minutos} mientras haya frescas
que pudrir. Casos borde a mencionar: sin frescas ($0$), fresca aislada
($-1$), todo vacío ($0$).
